\paragraph{}
Ett antal urvalskriterier togs fram för att nå ett relevant resultat från litteraturstudien.
Valet av databaser grundade sig i att det fanns relativt lite publicerad forskning kopplat till uppsatsens område.
Exempelvis saknade ACM:s databas relevant publicerad forskning och användes därför inte.
Sökträffarna från IEEE:s databas täckte forskningsområdet och har därför använts.
Sökning i Google scholar erhöll många sökträffar vilket utgjorde ett komplement till resultaten från IEEE:s databas.
Två söksträngar användes: \emph{(\lq driver advisory system\rq AND railway)} och \emph{(\lq driver advisory systems\rq AND railway)}.
\begin{enumerate}
    \item Inklusionskriterium
    \begin{enumerate}
        \item Skriven på svenska eller engelska
        \item Publicerad i vedertaget vetenskapligt forum
        \item Beskriver SFERA-standarden eller annat kommunikationsprotokoll mellan lokförare och tågklarerare
        \item Beskriver datat som tillhandahålles av TMS
    \end{enumerate}
    \item Exklusionskriterium
    \begin{enumerate}
        \item Publicerad innan 2010
        \item Berör utomeuropeisk järnväg
    \end{enumerate}
\end{enumerate}
\paragraph{}
Urvalsprocessen innehöll fyra steg.
Första steget innebar att utesluta verk baserat på titel.
Här sorterades eventuella dubbletter från annan databas bort.
För andra steget lästes abstract och för tredje steget lästes utöver abstract sammanfattning/conclusion (motsv.).
Kvarvarande artiklar lästes i sin helhet.
För samtliga steg prövades artiklarna mot inklusions- och exklusionskriterierna för eventuell uteslutning.
Antal sökträffar per steg redovisas nedan i tabell~\ref{tab:UrvalsprocessIEEE} för IEEE:s databas samt tabell~\ref{tab:UrvalsprocessGoogleScholar} för Google scholar:s databas.
Utöver kvarvarande artiklar som återstod efter litteraturstudiens urvalsprocess ingick \emph{n} interna dokument från Trafikverket.
Sammanlagt inkluderades \emph{n} verk i litteraturstudien.
\begin{table}[hbt!]
    \centering
    \begin{tabular}{|l|c|c|}
        \hline
        Urvalssteg & Antal träffar \\
        \hline
        0. Träffar & 17 \\
        1. Titel & 12 \\
        2. Abstract & 7 \\
        3. Slutsats & 4 \\
        4. Hela artikeln & 3 \\
        \hline
    \end{tabular}
    \caption{Urvalsprocess IEEE}
    \label{tab:UrvalsprocessIEEE}
\end{table}\begin{table}[hbt!]
    \centering
    \begin{tabular}{|l|c|}
        \hline
        Urvalssteg & Antal träffar \\
        \hline
        0. Träffar & 386 \\
        1. Titel & ? \\
        2. Abstract & ? \\
        3. Slutsats & ? \\
        4. Hela artikeln & ? \\
        \hline
    \end{tabular}
    \caption{Urvalsprocess Google scholar}
    \label{tab:UrvalsprocessGoogleScholar}
\end{table}